% This must be in the first 5 lines to tell arXiv to use pdfLaTeX, which is strongly recommended.
\pdfoutput=1
% In particular, the hyperref package requires pdfLaTeX in order to break URLs among lines.

\documentclass[11pt]{article}

% Change "review" to "final" to generate the final (sometimes called camera-ready) version.
% Change to "preprint" to generate a non-anonymous version with page numbers.
\usepackage[final]{acl}

% Standard package includes
\usepackage{times}
\usepackage{latexsym}

% For proper rendering and hyphenation of words containing Latin characters (including in bib files)
\usepackage[T1]{fontenc}
% For Vietnamese characters
% \usepackage[T5]{fontenc}
% See https://www.latex-project.org/help/documentation/encguide.pdf for other character sets

% This assumes your files are encoded as UTF8
\usepackage[utf8]{inputenc}

% This is not strictly necessary, and may be commented out,
% but it will improve the layout of the manuscript,
% and will typically save some space.
\usepackage{microtype}

% This is also not strictly necessary, and may be commented out.
% However, it will improve the aesthetics of text in
% the typewriter font.
\usepackage{inconsolata}

%Including images in your LaTeX document requires adding
%additional package(s)
\usepackage{graphicx}

% Add TikZ package for creating diagrams
\usepackage{tikz}
\usetikzlibrary{shapes.geometric, arrows.meta, positioning, fit, backgrounds}

% If the title and author information does not fit in the area allocated, uncomment the following
%
%\setlength\titlebox{<dim>}
%
% and set <dim> to something 5cm or larger.

\title{Shouth NLP at SemEval-2025 Task 3: Multilingual Fact-Checking Retrieval Using Contrastive Learning}
\author{}
\author{Santiago Lares Harbin \\
  Universidad de San Andres \\
  \texttt{slaresharbin@udesa.edu.ar}}

\begin{document}
\maketitle
\begin{abstract}
We present a multilingual fact-checking retrieval system for the SemEval-2025 task of matching social media posts with relevant fact checks. Our approach is based on a fine-tuned multilingual E5 model and in the experimentation with several prefixes and language-specific corpus filtering. The system achieves a Success@10 score of $0.867$ on the official test set, although with a big variance between languages. Our analysis discover interesting patterns in cross-lingual transfer, with strong results on Malaysian and Thai languages. We make our code public for further research and development.
\end{abstract}

\section{Introduction}



The proliferation of misinformation on social media has created an urgent need for automated fact-checking systems that can efficiently detect and verify false claims, given the immense volume of information shared online \cite{zhou2020fakenews}. Most automated fact-checking system consist of three steps \cite{guo2022survey}: first, the claim detection (is this check-worthy?); then, the evidence retrieval (given a claim, retrieve relevant information to verify it); and lastly and claim verification (given a claim and evidence, define if the claim is true or false). The SemEval-2025 Task \cite{sebi2025multilingual} addresses the second of these steps by requiring participants to develop evidence-retrieval systems for social media posts, out of a large corpus of $272,447$ fact checks among multiple languages.

In this paper, we present our approach to this task, which uses a fine-tuned multilingual embedding model based on the E5 architecture \cite{wang2024multilingual}. Our system employs a contrastive learning framework to optimize the semantic matching between posts and fact checks. We focus in particular on effective multilingual representation learning and cross-lingual retrieval functions. The code for our system is publicly available on GitHub\footnote{\url{https://github.com/sanlares/semeval_2025}}.

Our contributions are the following:
\begin{itemize}
  \item an analysis of various model architectures and their performance among different languages
  \item an experimentation of how input prefixes affect retrieval accuracy
  \item an exloration of the importance of language-specific corpus filtering
\end{itemize}

Our system achieved an average Success@10 score of 0.867 on the official test set, ranking 20th out of 28 published participants.

\section{Background}
%
% % Acá tenés que citar más papers...
% % Te tiro un párrafo de un texto no publicado, fijate si lo podés reutilizar acá
%
% %La tarea de fact-checking consiste en evaluar la veracidad de afirmaciones, como por ejemplo “La pobreza en Argentina es del 50%”. Algunos sitios como Snopes, Chequeado, y Full Fact dedican muchísimo esfuerzo humano a esta tarea, tanto en el contexto de afirmaciones de políticos como de desarmar noticias falsas o fake news [1]. Recientemente, y en concordancia con el avance del PLN, ha tenido lugar una creciente atención a la posibilidad de automatizar esta tarea. La automatización del fact-checking puede descomponerse en la concatenación de tres subtareas [16]: (i) la detección de afirmaciones; (ii) la recuperación de evidencia; y (iii) el veredicto sobre la afirmación. La detección de afirmaciones se basa también en evaluar la importancia de la aseveración y si es digna de ser evaluada (check-worthiness) y está relacionada a otras tareas como la detección de rumores [16]. La detección automática de fake news ha estado dominada hasta el momento por cuestiones de estilo de escritura y otros datos no factuales [1] mientras el fact-checking, en cambio, permanece como un caso elusivo a las actuales técnicas de LLMs. Si bien hay múltiples conjuntos de datos para las distintas etapas de fact-checking, mayormente están construidos en inglés. Podemos entre ellos mencionar a CheckThat21 [19] y las distintas versiones de FEVER [9], donde sólo el primero de estos tiene una porción de sus datos en español.
%
% [1] Zhou, X., & Zafarani, R. (2020). A survey of fake news: Fundamental theories, detection methods, and opportunities. ACM Computing Surveys (CSUR), 53(5), 1-40.
% [2] Zhang, Z., Han, X., Liu, Z., Jiang, X., Sun, M., & Liu, Q. (2019, July). ERNIE: Enhanced Language Representation with Informative Entities. In Proceedings of the 57th Annual Meeting of the Association for Computational Linguistics (pp. 1441-1451).
% [3] Wang, X., Gao, T., Zhu, Z., Zhang, Z., Liu, Z., Li, J., & Tang, J. (2021). KEPLER: A unified model for knowledge embedding and pre-trained language representation. Transactions of the Association for Computational Linguistics, 9, 176-194.
% [4] Pérez, J. M., Furman, D. A., Alemany, L. A., & Luque, F. (2021). RoBERTuito: a pre-trained language model for social media text in Spanish. arXiv preprint arXiv:2111.09453. Enviado a LREC 2022
% [5] Liu, W., Zhou, P., Zhao, Z., Wang, Z., Ju, Q., Deng, H., & Wang, P. (2020, April). K-bert: Enabling language representation with knowledge graph. In Proceedings of the AAAI Conference on Artificial Intelligence (Vol. 34, No. 03, pp. 2901-2908).
% [6] Arango, A., Pérez, J., & Poblete, B. (2019, July). Hate speech detection is not as easy as you may think: A closer look at model validation. In Proceedings of the 42nd international ACM SIGIR conference on research and development in information retrieval.
% [7] Devlin, J., Chang, M. W., Lee, K., & Toutanova, K. (2019, June). BERT: Pre-training of Deep Bidirectional Transformers for Language Understanding. In Proceedings of the 2019 Conference of the North American Chapter of the Association for Computational Linguistics: Human Language Technologies, Volume 1 (Long and Short Papers) (pp. 4171-4186).
% [8] Bordes, A., Usunier, N., Garcia-Duran, A., Weston, J., & Yakhnenko, O. (2013). Translating embeddings for modeling multi-relational data. Advances in neural information processing systems, 26.
% [9] Aly, R., Guo, Z., Schlichtkrull., et al (2021, June). FEVEROUS: Fact Extraction and VERification Over Unstructured and Structured information. In Thirty-fifth Conference on Neural Information Processing Systems Datasets and Benchmarks Track (Round 1).
% [10] Logan, R., Liu, N. F., Peters, M. E., Gardner, M., & Singh, S. (2019, July). Barack’s Wife Hillary: Using Knowledge Graphs for Fact-Aware Language Modeling. In Annual Meeting of the Association for Computational Linguistics (ACL).
% [11] Bender, E. M., Gebru, T., McMillan-Major, A., & Shmitchell, S. (2021, March). On the Dangers of Stochastic Parrots: Can Language Models Be Too Big?. In Proceedings of the 2021 ACM Conference on Fairness, Accountability, and Transparency (pp. 610-623).
% [12] Poletto, F., Basile, V., Sanguinetti, M., Bosco, C., & Patti, V. (2021). Resources and benchmark corpora for hate speech detection: a systematic review. Language Resources and Evaluation, 55(2), 477-523.
% [13] Chung, Y. L., Kuzmenko, E., Tekiroğlu, S. S., & Guerini, M. (2019, July). CONAN-COunter NArratives through Nichesourcing: a Multilingual Dataset of Responses to Fight Online Hate Speech. In Proceedings of the 57th Annual Meeting of the Association for Computational Linguistics (pp. 2819-2829).
% [14] Damián A. Furman, Pablo E. Torres, José A. Rodríguez, Laura Alonso Alemany, Diego Letzen, Maria Vanina Martinez. Parsimonious Argument Annotations for Hate Speech Counter-narratives. Enviado a LREC 2022
% [15] Alex Wang, Amanpreet Singh, Julian Michael, Felix Hill, Omer Levy, and Samuel Bowman. GLUE: A multi-task benchmark and analysis platform for NLU. In Proceedings of the 2018 EMNLP Workshop BlackboxNLP: Analyzing and Interpreting Neural Networks for NLP, pages 353–355, Brussels, Belgium, November 2018
% [16] Guo, Z., Schlichtkrull, M., & Vlachos, A. (2022). A survey on automated fact-checking. TACL, 10, 178-206.
% [17] Lundberg, S. M., & Lee, S. I. (2017). A unified approach to interpreting model predictions. NIPS, 30.
% [18] Choi, E., Levy, O., Choi, Y., & Zettlemoyer, L. (2018, July). Ultra-Fine Entity Typing. In Proceedings of the 56th Annual Meeting of the Association for Computational Linguistics (Volume 1: Long Papers) (pp. 87-96).
% [19] Nakov, P., Da San Martino, G., Elsayed, T., Barrón-Cedeno, A., Míguez, R., Shaar, S., ... & Mandl, T. (2021, March). The CLEF-2021 CheckThat! lab on detecting check-worthy claims, previously fact-checked claims, and fake news. In European Conference on Information Retrieval (pp. 639-649). Springer, Cham.
% [20] Schmidt, A., & Wiegand, M. (2019, January). A survey on hate speech detection using natural language processing. In Proceedings of the Fifth International Workshop on Natural Language Processing for Social Media, 2017, Valencia, Spain (pp. 1-10).
%

Automated fact-checking has emerged as a critical tool in combating the spread of misinformation online. The SemEval-2025 task focuses especially on multilingual fact-check retrieval, where the goal is to match social media posts with relevant fact-checks from a large multilingual corpus.

% Convertir a itemize
The task involves two main challenges: (1) semantic matching between posts and fact checks, which may be expressed differently despite covering the same claim, and (2) effective cross-lingual retrieval, where a post in one language must be matched with fact checks potentially written in different languages.

Recent advances in multilingual embedding models have shown promising results for cross-lingual retrieval tasks \cite{feng2022languageagnosticbertsentenceembedding, litschko2022crosslingual}. These models learn a shared semantic space between languages, allowing for effective comparison of text regardless of the source language. In particular, contrastive learning approaches have proven effective for training such models, as they explicitly optimize for similarity between related texts while pushing unrelated texts apart in the embedding space.

Our approach builds upon these advances, utilizing the multilingual E5 model architecture \cite{wang2024multilingual} as a foundation. This model has demonstrated strong performance on multilingual retrieval benchmarks by combining transformer-based encoders with effective contrastive pre-training. We additionally fine-tune this model on the task-specific data to optimize its performance for fact-check retrieval.

\section{System Overview}

Our system addresses the challenge of retrieving relevant fact checks from a large corpus of 272,447 documents by implementing a fine-tuned contrastive learning approach based on the multilingual E5 model architecture. The system consists of three main components: (1) a text embedding model, (2) a contrastive learning framework, and (3) a retrieval pipeline.

\subsection{Model Architecture}

% Acá citá el paper de E5! y describí brevemente lo que hace

The core of our system is built upon the Multilingual E5 model \cite{wang2024multilingual} which we improve with a custom architecture designed for fact-checking retrieval. The model consists of:

% Estaría buenísimo convertir esto en un gráfico...
\begin{itemize}
    \item A transformer encoder (base E5 model)
    \item A mean pooling layer for sentence representation
    \item A dense projection layer with GELU activation
    \item L2 normalization for cosine similarity optimization
\end{itemize}

The model processes both posts and fact checks using language-specific prefixes (\texttt{query:} and \texttt{passage:} respectively) to optimize the semantic matching between them.

\subsection{Training}

% citar esta pérdida...
We trained our model using a the Multiple Negatives Ranking Loss (MNRL), which can be formulated as:

\begin{equation}
L = -\log \frac{\exp(sim(x_i, x_i^+)/\tau)}{\sum_{j=1}^{N} \exp(sim(x_i, x_j)/\tau)}
\end{equation}

where:
\begin{itemize}
    \item $x_i$ represents the anchor text embedding (post)
    \item $x_i^+$ represents the positive pair embedding (matching fact-check)
    \item $x_j$ represents all other samples in the batch (negative pairs)
    \item $\tau$ is a temperature parameter that controls the sharpness of the distribution
    \item $sim(\cdot,\cdot)$ is the cosine similarity function
\end{itemize}

We used in-batch negatives for computational efficiency while maintaining effective contrastive learning. The temperature parameter $\tau$ is dynamically adjusted during training to optimize the learning process.

\subsection{Retrieval Pipeline}

The retrieval process follows these steps:

\begin{enumerate}
    \item \textbf{Language Detection}: The system first identifies the language of the input post.
    \item \textbf{Corpus Filtering}: The fact-check corpus is filtered to prioritize documents in the same language as the input post.
    \item \textbf{Embedding Generation}: The model generates embeddings for both the input post and the filtered fact checks.
    \item \textbf{Similarity Ranking}: The system computes Pearson cosine similarity between the post and fact check embeddings.
    \item \textbf{Top-K Selection}: The 10 most similar fact checks are retrieved and ranked.
\end{enumerate}

This pipeline is optimized for both monolingual and cross-lingual retrieval, with special consideration for language-specific distinctions in the embedding space.

\section{Experimental Setup}

\subsection{Dataset}

% JM: ¿vos lo partiste o ya venía así?

The dataset used for this study was derived from the official SemEval-2025 competition data. We divided the training set into training (80\%) and validation (20\%) subsets while maintaining the language distribution for both sets.


For preprocessing, we used the script provided by the competition organizers.For posts, we concatenated the OCR-extracted text (\texttt{ocr} column) with the manual transcription (\texttt{text} column) from the \texttt{posts.csv} file. Similarly, for fact checks, we combined the \texttt{claim} and \texttt{title} columns from the \texttt{fact\_checks.csv} file.

In addition to fine-tuning our primary model, we conducted extensive experiments with various base models available in the SentenceTransformers framework, testing different prefix combinations that yielded varying performance results.

\subsection{Implementation Details}


We used the \textit{multilingual-e5-large-instruct} found at the Huggingface model hub \cite{wang2024multilingual} as our base model. We fine-tuned the retrieval model using the SentenceTransformers library \cite{reimers-gurevych-2019-sentence}, and used PyTorch \cite{paszke2019pytorch} to implement our models and losses.

The model was trained with the following hyperparameters:

\begin{table}[h]
\centering
\small
\begin{tabular}{ll}
\hline
\textbf{Parameter} & \textbf{Value} \\
\hline
Maximum Sequence Length & 512 tokens \\
Pooling Mode & Mean \\
Embedding Dimension & 384 \\
Batch Size & 4 \\
Evaluation Batch Size & 16 \\
Gradient Accumulation Steps & 2 \\
Learning Rate & 2e-5 \\
Temperature & 0.05 \\
Mixed Precision Training & Enabled \\
\hline
\end{tabular}
\caption{Model training hyperparameters}
\label{tab:hyperparams}
\end{table}

\subsection{Evaluation}

The primary evaluation metric for our system was Success@10 (S@10), which measures the proportion of posts for which the correct fact check appears in the top 10 retrieved results. We evaluated the model's performance in both monolingual and cross-lingual settings:

\begin{itemize}
    \item \textbf{Monolingual}: Retrieval within the same language
    \item \textbf{Cross-lingual}: Retrieval across different languages
\end{itemize}

During training, we also monitored:
\begin{itemize}
    \item Cosine similarity between positive pairs
    \item Cosine similarity between negative pairs
    \item Training loss convergence
    \item Validation performance at regular intervals
\end{itemize}

\section{Results}

We evaluate our system's performance on both the validation and test sets, analyzing the impact of different model configurations and language-specific performance. Our system achieved an average Success@10 score of \textbf{0.867} on the official test set, ranking 20th out of 28 published participants.

\subsection{Official Test Results}

\begin{table}[t]
  \centering
  \small
  \begin{tabular}{lcc}
    \hline
    \textbf{Language} & \textbf{Success@10} & \textbf{Deviation from avg.} \\
    \hline
    Malaysian (msa) & \textbf{0.978} & +12.8\% \\
    Thai (tha) & 0.940 & +8.4\% \\
    Arabic (ara) & 0.912 & +5.2\% \\
    French (fra) & 0.894 & +3.1\% \\
    Spanish (spa) & 0.866 & -0.1\% \\
    German (deu) & 0.858 & -1.0\% \\
    Turkish (tur) & 0.828 & -4.5\% \\
    Polish (pol) & 0.806 & -7.0\% \\
    English (eng) & 0.796 & -8.2\% \\
    Portuguese (por) & 0.796 & -8.2\% \\
    \hline
    \textbf{Average} & 0.867 & -- \\
    \hline
  \end{tabular}
  \caption{Official test set results of the fine-tuned multilingual-e5-large-instruct model by language. Languages are ordered by Success@10 score. Best performance is shown in \textbf{bold}. The model shows strong performance on Malaysian and Thai languages.}
  \label{tab:test-results-by-language}
\end{table}

\vspace{1cm}

The results demonstrate considerable variability in model performance amid languages, with Malaysian (msa) showing the strongest performance at 0.978 and English (eng) and Portuguese (por) showing the lowest at 0.796. This disparity suggests that linguistic factors may play a significant role in retrieval performance.

\subsection{Model Configuration Analysis}

Table~\ref{tab:base-models} shows the performance comparison of different base models and input prefix configurations on the validation set. The results demonstrate that the multilingual-e5-large-instruct model with appropriate prefixes consistently outperforms other configurations.

\begin{table*}[t]
  \centering
  \small
  \begin{tabular}{p{5cm}p{3cm}p{3cm}c}
    \hline
    \textbf{Base Model} & \textbf{Posts Prefix} & \textbf{Facts Prefix} & \textbf{Success@10} \\
    \hline
    \texttt{multilingual-e5-large-instruct} \cite{wang2024multilingual} & passage: & query: & \textbf{0.834} \\
    \texttt{multilingual-e5-large-instruct} \cite{wang2024multilingual} & -- & -- & 0.821 \\
    \texttt{multilingual-e5-small} \cite{wang2024multilingual} & -- & -- & 0.795 \\
    \texttt{e5-large-v2} \cite{wang2022text} & -- & -- & 0.758 \\
    \texttt{modernbert-embed-base} \cite{nussbaum2024nomic} & \texttt{search\_query:} & \texttt{search\_document:} & 0.779 \\
    \texttt{paraphrase-multilingual-mpnet-base-v2} \cite{reimers2019sentence} & -- & -- & 0.736 \\
    \hline
  \end{tabular}
  \caption{Performance comparison of base models on the validation set. Models are ordered by Success@10 score. Best result is shown in \textbf{bold}.}
  \label{tab:base-models}
\end{table*}

\subsection{Impact of Input Prefixes}

Our experiments revealed that the choice of input prefixes significantly impacts model performance. Using the recommended prefixes on multilingual-e5-large-instruct base model ('passage:' and 'query:') improved the average Success@10 score by 0.013 points (from 0.821 to 0.834) compared to using no prefixes. This improvement was consistent across all languages, with the largest gains observed in:

\begin{itemize}
    \item English: +0.022 (from 0.766 to 0.788)
    \item Spanish: +0.012 (from 0.857 to 0.869)
    \item Portuguese: +0.018 (from 0.781 to 0.799)
\end{itemize}

These results suggest that appropriate input formatting plays a vital role in optimizing the model's cross-lingual understanding and retrieval functions.

\subsection{Corpus Reduction Analysis}

Our experiments showed that performance improves significantly when the fact-check corpus is reduced to include only documents in the same language as the query post. For example, when retrieving from a language-specific subset of approximately 16,000 fact checks (versus the full corpus of 272,447), Success@10 scores improved by an average of 10.2\%. For more detailed experiments on corpus reduction with the fine-tuned model, see Table~\ref{tab:prefix-experiments} in Appendix~\ref{sec:appendix-results}.

\section{Conclusion}

In this paper, we presented a multilingual fact-checking retrieval system based on fine-tuned contrastive learning with the E5 model architecture. Our approach achieved an average Success@10 score of 0.867 on the official SemEval-2025 test set, with specifically strong performance on non-Latin script languages such as Malaysian and Thai.

Our analysis revealed several key outcomes:

\begin{itemize}
    \item The base model multilingual-e5-large-instruct consistently outperforms other options.
    \item Appropriate input prefixes can provide substantial performance gains, improving Success@10 by 0.013 points on average.
    \item There is considerable variability in model performance between languages, suggesting that language-specific tuning could yield extended improvements.
    \item Corpus reduction by language dramatically improves retrieval accuracy, highlighting the importance of effective pre-filtering strategies.
\end{itemize}

Subsequent work could focus on addressing the performance discrepancy among languages, perhaps through language-specific fine-tuning or more sophisticated cross-lingual transfer techniques. Moreover, exploring ensemble methods that combine multiple embedding models might improve retrieval accuracy for challenging languages.

\section{Acknowledgments}

We would like to thank the SemEval-2025 task organizers for creating this challenging and impactful task, and for providing thorough evaluation metrics and datasets. We also acknowledge the support of Juan Manuel Pérez from Universidad de San Andres.

% Bibliography entries for the entire Anthology, followed by custom entries
%\bibliography{anthology,custom}
% Custom bibliography entries only
\bibliography{anthology, custom}

\appendix

\section{Detailed Experimental Results}
\label{sec:appendix-results}

This appendix presents supplementary experimental results and analyses that complement the main outcomes discussed in the paper.

\subsection{Corpus Reduction Analysis}

Table~\ref{tab:prefix-experiments} shows the impact of corpus reduction on the fine-tuned model's performance. The results demonstrate that using language-specific subsets of the fact-check corpus significantly improves retrieval performance.

\begin{table}[h]
  \centering
  \small
  \begin{tabular}{p{2.5cm}p{1.8cm}p{1.8cm}c}
    \hline
    \textbf{Model} & \textbf{Posts Prefix} & \textbf{Facts Prefix} & \textbf{S@10} \\
    \hline
    multilingual-e5 & query: & passage: & \textbf{0.934} \\
    multilingual-e5 & -- & -- & 0.933 \\
    multilingual-e5 & post: & fact check: & 0.934 \\
    multilingual-e5 & This is a post... & This is a fact... & 0.936 \\
    multilingual-e5 & <[language]> & -- & 0.933 \\
    multilingual-e5 & The following... & The following... & 0.934 \\
    \hline
  \end{tabular}
  \caption{Fine-tuned model performance with different configurations when using a reduced corpus. All experiments were conducted on the validation set with a corpus containing only fact checks in the same language as the query posts, resulting in higher overall scores compared to the full corpus evaluation.}
  \label{tab:prefix-experiments}
\end{table}

Our experiments showed that performance improves significantly when the fact-check corpus is reduced to include only documents in the same language as the query post. For example, when retrieving from a language-specific subset of approximately 16,000 fact checks (versus the full corpus of 272,447), Success@10 scores improved by an average of 10.2\%.

\subsection{Fine-tuning Impact}

The base model `intfloat/multilingual-e5-large-instruct` \cite{wang2024multilingual} improved from an average Success@10 of 0.834 to 0.867 when fine-tuned with the training dataset, along with the recommended query and passage prefixes. This improvement was consistent for all languages tested in the competition.

\end{document}
